\slajdid{09_2-s045}
\begin{frame}[label=09_2-s045]
odnosno za zadato $V_S$ odnosi treba da budu
\[\frac{W_p}{W_n}=\frac{\mu_nV_{DSnsat}\left(V_S-V_{Tn}-\frac{V_{DSnsat}}2\right)}{\mu_pV_{DSpsat}\left(V_S-V_{DD}-V_{Tp}-\frac{V_{DSpsat}}2\right)}\]
U slučaju da tačku $V_S$ podešavamo na polovinu napona napajanja i uz uslov $v_{satn}=-v_{satp}$
\[\frac{W_p}{W_n}=\frac{\frac{V_{DD}}2-V_{Tn}-\frac{V_{DSnsat}}2}{\frac{V_{DD}}2+V_{Tp}+\frac{V_{DSpsat}}2}\]
Što je identično i prethodnom izrazu i takođe se kod dugog kanala $V_{DSnsat}>\frac{V_{DD}}2-V_{Tn}$ svodi na
\[\frac{W_p}{W_n}=\frac{V_{DSnsat}}{|V_{DSpsat}|}=\frac{\mu_n}{\mu_p}\]
\end{frame}
